\documentclass[10pt]{mypackage}

\usepackage[uprightgreeks,light]{kpfonts}
\renewcommand{\coloneq}{\coloneqq}
\usepackage{homework}
%\usepackage{notes}

%\usepackage[ backend=bibtex, style = alphabetic, sorting=ynt ]{biblatex}
%\addbibresource{  }

\usepackage{parskip}

\fancyhf{}
\fancyhead[R]{Avinash Iyer}
\fancyhead[L]{Algebraic Topology: Homework 22}
\fancyfoot[C]{\thepage}

\setcounter{secnumdepth}{0}

\begin{document}
\RaggedRight
\begin{problem}[Problem 1]
  Prove that $ \widetilde{H}_q(X) = H_q(X) $ for $q\neq 0$ and $H_0\left( X \right) \cong \widetilde{H}_0(X) \oplus \Z$ for any space $X$.
\end{problem}
\begin{solution}
  For any $q\neq 0$, we observe that $\widetilde{C}_q(X) = C_q(X)$, and since the maps $\widetilde{\partial}_{q}$ for any $q > 1$ on the reduced singular chain complex are identical to $\partial_q$ on the original singular chain complex. In particular, this means that the corresponding boundaries $ \mathcal{B}_q $ and $\mathcal{Z}_q$ are identical for any $q \geq 1$, so $ \widetilde{H}_q(X) = H_q(X) $ for all $q \neq 0$.

  When it comes to $\widetilde{H}_0$, we observe that any homology class $\left[ z \right]\in \widetilde{H}_0(X)$ is contained in $H_0(X)$ as $\ker\left( \ve \right)\leq C_0\left( X \right)\cong H_0\left( X \right)$. Furthermore, the augmentation homomorphism $\ve\colon C_0\left( X \right)\rightarrow \Z$ is surjective. This gives the following short exact sequence, which splits as $C_0(X)$ is a free abelian group/$\Z$-module.
  \begin{center}
    % https://tikzcd.yichuanshen.de/#N4Igdg9gJgpgziAXAbVABwnAlgFyxMJZABgBpiBdUkANwEMAbAVxiRGJAF9T1Nd9CKAIzkqtRizYAdKQHcssPA1jAAEpwD6xABQANAJRceIDNjwEiAJlHV6zVohCqtew915mBRAMw3x96SkAWzocAAsAIwjgAC1OIw9+CxQAFj87SUcOTjEYKABzeCJQADMAJwggpDIQHAgkIXcQcsqG6jqka39MkBl8HDoE5oqqxC6OxF9uh16pejKYNGwGCyaW0amJlJzOIA
    \begin{tikzcd}
    0 \arrow[r] & \widetilde{C}_0(X) \arrow[r, "\iota"] & C_0(X) \arrow[r, "\varepsilon"] & \mathbb{Z} \arrow[r] & 0
    \end{tikzcd}
  \end{center}
  This gives that $C_0\left( X \right)\cong \mathcal{C}_0\left( X \right)\oplus \Z$. Since $B_0(X)\cong B_0\left( X \right)\oplus \set{0}$, and $B_0\left( X \right) = \img\left( \partial_1 \right)$ in both reduced and original singular homology, the rules of direct sums of modules gives
  \begin{align*}
    H_0\left( X \right) &= C_0\left( X \right)/B_0\left( X \right)\\
                        &\cong \left( \widetilde{C}_0(X)\oplus \Z \right)/\left( B_0(X)\oplus \set{0} \right)\\
                        &\cong \widetilde{C}_0(X)/B_0(X)\oplus \Z\\
                        &= \widetilde{H}_0(X)\oplus \Z.
  \end{align*}
\end{solution}
\begin{problem}[Problem 2]
  For any space $X$, let $SX$ denote the suspension of $X$. Prove that for any $q\geq 0$, there is an isomorphism
  \begin{align*}
    \widetilde{H}_q(X) &\cong \widetilde{H}_{q+1}\left( SX \right).
  \end{align*}
\end{problem}
\begin{solution}
  We view the suspension $SX$ as $X\times \left[ -1,1 \right]/\left( X\times \set{-1} \right)/\left( X\times \set{1} \right)$, for the sake of simplicity. Consider the subsets
  \begin{align*}
    U &= \set{\left[ \left( x,t \right) \right] | x\in X,t \leq -1/2}\\
    A &= \set{\left[ \left( x,t \right) \right] | x\in X, t \leq 0}.
  \end{align*}
  First, observe that $A\cong CX$, which is a contractible space, so $\widetilde{H}_q(A) = 0$ for all $q$, meaning that $ \widetilde{H}_q\left( SX \right)\cong \widetilde{H}_q\left( SX,A \right) $ as can be seen by the long exact sequence.
  \begin{center}
    % https://tikzcd.yichuanshen.de/#N4Igdg9gJgpgziAXAbVABwnAlgFyxMJZARgBoAGAXVJADcBDAGwFcYkQAdDgdy1j0axgACQC+AfQCOACgDKADQCUIUaXSZc+QinIVqdJq3Zde-LIJgiJMgILLV67HgJEATHpoMWbRJx58YASExKTl5UjsVNRAMJy0iAGYPA29jfzMLK3FgSQBaYlFpSNF9GCgAc3giUAAzACcIAFskMhAcCCRyBxB6ps6adqR3FKNfLgBjAnKo2obmxGHBxASS0SA
\begin{tikzcd}
\widetilde{H}_q(A) \arrow[r] & \widetilde{H}_q(SX) \arrow[r, "\cong"] & {\widetilde{H}_q(SX,A)} \arrow[r] & \widetilde{H}_{q-1}(A)
\end{tikzcd}
  \end{center}
  Since $ \overline{U}\subseteq A^{\circ}$, it follows from excision that $\widetilde{H}_{q}\left( SX,A \right) = \widetilde{H}_{q}\left( SX\setminus U,A\setminus U \right)$. Observe now that by a deformation retraction, we have that $ A\setminus U\simeq X $, while $SX\setminus U\simeq CX$, where $CX$. Since $CX$ is contractible, the long exact sequence on $ H_{\ast}\left( CX,X \right) $ gives that $\widetilde{H}_{q+1}\left( CX,X \right)\cong \widetilde{H}_{q}(X)$
  \begin{center}
    % https://tikzcd.yichuanshen.de/#N4Igdg9gJgpgziAXAbVABwnAlgFyxMJZABgBpiBdUkANwEMAbAVxiRAB12B3LWPB2MAASAXwD6wAI4BqAIwiAFAGEAGgEoQI0uky58hFLPJVajFm048+WATGHipcxatLrN2kBmx4CRAEzG1PTMrIgc3Lww-IKiEpKKblo63vpEAMyBpiEWEda29nHOiSYwUADm8ESgAGYAThAAtkhkIDgQSPIedY0d1G1IAVnmYZwAxgRl7jX1TYiD-YhpIhQiQA
    \begin{tikzcd}
    \widetilde{H}_{q+1}(CX) \arrow[r] & {\widetilde{H}_{q+1}(CX,X)} \arrow[r, "\cong"] & \widetilde{H}_{q}(X) \arrow[r] & \widetilde{H}_{q}(CX)
    \end{tikzcd}
  \end{center}
  Combined, we have the following chain of isomorphisms
  \begin{align*}
    \widetilde{H}_{q+1}\left( SX \right) &\cong \widetilde{H}_{q+1} \left( SX,A \right) \tag{i}\\
                                         &\cong \widetilde{H}_{q+1} \left( SX\setminus U,A\setminus U \right)\tag{ii}\\
                                         &\cong \widetilde{H}_{q+1}\left( CX,X \right) \tag{iii}\\
                                         &\cong \widetilde{H}_{q}\left( X \right) \tag{iv}
  \end{align*}
  where (i) was from the long exact sequence on $ \widetilde{H}_{\ast}\left( SX,A \right) $, (ii) from excision, (iii) from homotopy invariance of reduced homology, and (iv) from the long exact sequence on $\widetilde{H}_{\ast}\left( CX,X \right)$.
\end{solution}
\end{document}
